\section{Cuts moderately larger than the girth}
\label{sec:four-high-connectivity-near-three-halves}

The preceding entropy argument tolerates slightly more ports than the girth.
This extension is separate from the main minimum-cut result.
All logarithms have base two.

\begin{lemma}
\label{lem:four-near-three-halves-rank}
Fix $0<\delta<3/2$.
For sufficiently large $g$, depending only on $\delta$, every normalized $k$-pole of rank $h\ge2$ satisfying
\[
 \operatorname{girth}(N)\ge g,
 \qquad k\le(3/2-\delta)g
\]
has
\[
 h\ge2^{g/2-5}.
\]
\end{lemma}

\begin{proof}
The theta-or-two-cycles argument gives
\[
 g\le2h+\frac{2k}{3}-2,
 \qquad h-1\ge\frac{\delta g}{3}.
\]
Hence the stationary non-backtracking entropy rate satisfies
\[
 \sigma=\frac{3(h-1)}{3(h-1)+k}
 \ge\frac{\delta}{\delta+3/2}.
\]
If $h\le2g^2$, then $n=2h+k-2\le6g^2$.
Lemma~\ref{lem:four-entropy-moore} would give
\[
 g\le4+2\left(1+\frac{3}{2\delta}\right)\log(6g^2),
\]
which fails for sufficiently large $g$.
Thus $h>2g^2$, and in particular $h-1\ge g^2$.

If $h\ge2^{g/2}$, the conclusion already holds.
Otherwise $n\le2^{g/2+2}$ and
\[
 \frac{2k\log n}{3(h-1)}
 \le\frac{3g(g/2+2)}{3g^2}<1
\]
for $g>4$.
The entropy bound gives $g<5+2\log n$ and hence $n>2^{(g-5)/2}$.
For sufficiently large $g$, this lower bound exceeds $3g$.
Since $k\le3g/2$, it follows that
\[
 h=\frac{n-k+2}{2}>\frac n4>2^{g/2-5}.
\]
\end{proof}

\begin{corollary}
\label{cor:four-near-three-halves-forests}
For each fixed $0<\delta<3/2$, there is a constant $C_\delta$ such that every normalized $k$-pole of rank $h\ge2$ satisfying
\[
 2\le k\le(3/2-\delta)g,
 \qquad \operatorname{girth}(N)\ge g,
\]
obeys, for every even port subset $T$,
\[
 \frac{F_T(N)}{2^h}
 \le\frac12+C_\delta g^{13/2}2^{-g/4}.
\]
\end{corollary}

\begin{proof}
Repeat the proof of Theorem~\ref{thm:path-high-girth-terminal-forests}.
Use Lemma~\ref{lem:four-near-three-halves-rank} for its rank lower bound.
The radius $s=\lfloor g/2-\log k-10\rfloor$ still gives a port-free cubic ball.
The resulting layer lower bound is
\[
 \frac{2^{g/2}}{2^{11}kg}
 \ge\frac{2^{g/2}}{3\cdot2^{10}g^2}.
\]
The layer upper bound has the same form, with $(k+1)^2\le(3g/2+1)^2$.
Increasing the constant in the excess bound absorbs these fixed factors.
The rank threshold depends only on $\delta$.
Increase $C_\delta$ once more to cover the remaining finite range of $g$.
\end{proof}

\begin{corollary}[A quarter-space bound for moderate cyclic cuts]
\label{cor:four-near-three-halves-cuts}
Fix $0<\delta<3/2$.
Let connected cubic graphs have girths $g\to\infty$.
Suppose a cut of order $k\le(3/2-\delta)g$ has two connected shores, each of rank at least two.
Then
\[
 c(G)\le\left(\frac14+o_\delta(1)\right)2^{\beta(G)}.
\]
The error is uniform over the graphs and the chosen cuts.
\end{corollary}

\begin{proof}
First make the ports distinct if necessary.
A vertex with two cut incidences is a leaf in its connected shore.
Move it to the other shore.
Its old shore remains connected with the same rank, while the other shore remains connected and gains rank one.
The cut order decreases by one.
Neither rank falls below two.
A vertex with three cut incidences would be isolated in its shore, which is impossible.
Repeat until both shores have minimum internal degree two.
Write $k'\le k$ for the resulting cut order and $h,j\ge2$ for the shore ranks.
They are normalized $k'$-poles and
\[
 r=\beta(G)=h+j+k'-1.
\]

If $k'=1$, no cycle crosses the cut.
The entropy rank bound makes both shore ranks tend to infinity, and the conclusion follows from $c(G)\le2^h+2^j$.
Assume $k'\ge2$.
Lemma~\ref{lem:four-near-three-halves-rank} makes both shore ranks at least $2^{g/2-5}$ for sufficiently large $g$.
Put $\varepsilon_g=C_\delta g^{13/2}2^{-g/4}=o(1)$.
For each nonempty even boundary subset, both shore forest counts are at most their affine-space sizes times $1/2+\varepsilon_g$.
Every crossing cycle restricts to one such forest in each shore.
Therefore the crossing-cycle count is at most
\[
 (2^{k'-1}-1)(1/2+\varepsilon_g)^2 2^{h+j}
 \le(1/2+\varepsilon_g)^2 2^r.
\]
The internal cycles contribute at most $2^h+2^j=o(2^r)$ uniformly.
This proves the result.
\end{proof}
